\input{_preamble}
\begin{document}
% In each pair the leading word is identical, so the sentinel that follows
% it must land at exactly the same x -- unless the judgment mark wrongly
% consumed horizontal space in the text block.
\ex. *Text JMAIN rest.

\ex. Text PMAIN rest.

\ex.
\a. ??\#Text JSUB rest.
\b. Text PSUB rest.

\ex.
\a. Level,
\a. *Text JROMAN rest.
\b. Text PROMAN rest.

\ex. \jdg{\dag}Text JMANUAL rest.

\ex. Text PMANUAL rest.

% A label typed at the head of an example -- which is where a label
% actually gets typed -- must not stop the judgment scan.  \label makes
% no ink, so the mark is still the first thing on the line and still
% belongs in the gutter; before it was skipped past, the marks stayed in
% the body and indented exactly the examples that are cross-referenced.
% The pairs below are the same displacement test as the ones above, and
% the \ref line is the other half of it: the label has to survive being
% carried across the mark, still name this example, and still be a label
% of the right KIND (\sublabel's letter is what \prefrange reads).
\ex. \label{j:main}*Text JLABEL rest.

\ex. Text PLABEL rest.

% ... and a label with no mark behind it must not displace anything
% either: skipping past it means the scan, not \ignorespaces, eats the
% space that follows it, so QLABEL is where a spurious one would show.
\ex. \label{j:nomark}Text QLABEL rest.

\ex.
\a. \sublabel{j:sublab}??\#Text JSUBLAB rest.
\b. Text PSUBLAB rest.

\exg. \label{j:gloss}*Text JGLOSS rest.\\
Wort GLOSSTIER Rest.\\

\exg. Text PGLOSS rest.\\
Wort PGLOSSTIER Rest.\\

REFCHECK \ref{j:main} \ref{j:sublab} \ref{j:gloss} \prefrange{j:main}{j:sublab}
\end{document}
